\bookchapter{פרק 10: תשלומים ניידים ושירותי חיים}

\chapterintro{האמת, הדור של ההורים שלי בהתחלה התנגד ל“לשלם בטלפון”. “הכסף בטלפון, מה אם מישהו יגנוב אותו?” הדאגה הזו טבעית, אבל המציאות היא: הסיכוי שתאבד מזומן מהכיס הרבה יותר גבוה מהסיכוי שתשלום בטלפון ייגנב. תשלום בטלפון זה לא לשים כסף בתוך הטלפון, זה רק מפתח, מפתח שפותח את חשבון הבנק שלך.}

\booksection{איך בדיוק לקשר כרטיס בנק?}

הזן את מספר הכרטיס. עכשיו רוב האפליקציות תומכות בזיהוי צילום של מספר הכרטיס, לא צריך להקליד ספרה ספרה.

\begin{enumerate}
  \item הבנק ישלח קוד אימות ב-SMS למספר הטלפון השמור שלך. הזן אותו, הגדר סיסמת תשלום בת 6 ספרות.
  \item סיסמת התשלום הזו, סיסמת כרטיס הבנק וסיסמת נעילת המסך – עדיף שתהיינה שלוש סיסמאות שונות, אל תגדיר את כולן זהות מטעמי נוחות.
  \item טיפ קטן: אם אתה מודאג מאבטחה, אפשר לפתוח חשבון אלקטרוני מסוג II או III באפליקציית הבנק, במיוחד לקישור לתשלום בטלפון. שים בכרטיס הזה רק כסף להוצאות שוטפות, גם אם יקרה משהו, ההפסד מוגבל.
  \item 2. קוד תשלום: "מי שרואה – גובה כסף"
\end{enumerate}

הרבה אנשים לא מבינים את ההבדל בין קוד תשלום לקוד קבלה, וההבדל הזה חשוב מאוד.

\booksection{קוד תשלום: “מי שרואה – גובה כסף”}

: הברקוד + קוד QR שאתה מציג לקופאי לסריקה. הקוד הזה שווה ערך ל“ארנק האלקטרוני” שלך, מי שסורק אותו יכול למשוך כסף מהחשבון שלך. לכן

\begin{itemize}
  \item \textbf{לעולם אל תצלם מסך ותשלח לאף אחד}. הקוד הזה מתרענן אוטומטית כל דקה, בדיוק כדי למנוע גניבה.\textbf{קוד קבלה}: אחרים סורקים את הקוד הזה כדי להעביר לך כסף. את הקוד הזה אפשר לצלם מסך ולשמור, אפשר לשלוח לאחרים, בטוח.
  \item \textbf{תכיר את ההבדל הזה, ותוכל להימנע מהרבה מלכודות הונאה. כל מי שמבקש ממך "לצלם מסך של קוד התשלום ולשלוח", במאה אחוז רמאי.}3. הגדרות אבטחה אלה – 5 דקות שיחסכו לך הרבה דאגות
\end{itemize}

תשלום בטביעת אצבע/זיהוי פנים

\booksection{הגדרות אבטחה אלה – 5 דקות שיחסכו לך הרבה דאגות}

\begin{itemize}
  \item \textbf{נעילת אבטחה}: גם ל-Alipay וגם ל-WeChat יש פונקציית "נעילת אבטחה". לאחר הפעלה, הכניסה לדף "ארנק" באפליקציה דורשת אימות זהות נוסף. נתיב: Alipay שלי הגדרות חשבון ואבטחה מרכז אבטחה נעילת אבטחה.
  \item \textbf{בדיקת תשלום ללא סיסמה}: הרבה אנשים חתמו על כל מיני תשלומים ללא סיסמה ושכחו מזה. לך ל"הגדרות תשלום" ב-Alipay ו-WeChat, מצא "תשלום ללא סיסמה/חיוב אוטומטי", ראה אילו שירותים חתומים, וסגור את המיותרים.
  \item \textbf{ניהול מכשירים מחוברים}: בדוק מדי פעם אילו מכשירים התחברו לחשבון שלך. אם אתה מוצא מכשיר לא מוכר, נתק מיד. נתיב: Alipay הגדרות חשבון ואבטחה ניהול מכשירים מחוברים.
  \item \textbf{4. אוטובוס ורכבת תחתית: הטלפון הוא כרטיס הנסיעה}כיום, ביציאה לנסיעה באוטובוס או רכבת תחתית, מעט אנשים משתמשים בכרטיס נסיעה פיזי. לסריקת הטלפון בשער, יש שתי אפשרויות:
\end{itemize}

\booksection{אוטובוס ורכבת תחתית: הטלפון הוא כרטיס הנסיעה}

: פתח כרטיס נסיעה וירטואלי ישירות בארנק Apple, ארנק Huawei, ארנק Xiaomi. היתרון של הפתרון הזה: בסריקה אין צורך לבטל את נעילת הטלפון, פשוט מקרבים, מהירות התגובה הגבוהה ביותר, ועובד גם בלי רשת. בחלק מהערים, לכרטיס הנסיעה הווירטואלי יש גם הנחות מעבר כמו לכרטיס הפיזי.

\begin{itemize}
  \item \textbf{קוד נסיעה של Alipay/WeChat}: חפש "קוד נסיעה" באפליקציה, פתח את קוד הנסיעה לאוטובוס או רכבת תחתית בעיר שלך. היתרון של הפתרון הזה: אין צורך בטעינה, מחויב ישירות מהיתרה או מכרטיס הבנק המקושר. אבל דורש חיבור לרשת, וצריך לפתוח את האפליקציה ולמצוא את הקוד לפני הסריקה.
  \item \textbf{אני אישית ממליץ יותר על כרטיס הנסיעה הווירטואלי בארנק הנייד. למרות שבפעם הראשונה צריך לטעון 10-20 שקל, אבל השימוש באמת הרבה יותר נוח, במיוחד כשרצים לרכבת התחתית, לא צריך לעצור ולבטל נעילה ולחפש אפליקציה.}5. תשלומי חשבונות, משלוחי אוכל, הזמנת מוניות
\end{itemize}

תשלומי חשמל, מים וגז

\booksection{תשלומי חשבונות, משלוחי אוכל, הזמנת מוניות}

\begin{itemize}
  \item \textbf{השוואת מחירי משלוחים}: אותה חנות ב-Meituan וב-Ele.me לעיתים קרובות מציעה מחירים ומבצעים שונים. לפני ההזמנה, תסתכל בשתיהן, לפעמים אפשר לחסוך כמה שקלים. בנוסף, פונקציית "הזמנה משותפת" של Meituan, להזמין יחד עם חברים, המחיר יכול להיות זול בהרבה.
  \item \textbf{בטיחות בנסיעות}: ב-Didi או Gaode, זכור להגדיר "איש קשר לחירום" ו"שיתוף מסלול". בנסיעה, האפליקציה תשתף אוטומטית את המסלול שלך עם איש הקשר שהגדרת. בנוסף, בנסיעות מאוחרות בלילה העדף "מונית" על פני "נסיעה מהירה", למוניות יש פיקוח ענפי, מקדם הבטיחות גבוה יותר.
  \item \textbf{6. ביטוח רפואי, ביטוח לאומי ושירותים ממשלתיים בטלפון}השינויים בשנים האחרונות גדולים מאוד. הרבה דברים שבעבר דרשו המתנה בתור באולם השירות, כיום אפשר לסיים תוך דקות ספורות בטלפון.
\end{itemize}

\booksection{ביטוח רפואי, ביטוח לאומי ושירותים ממשלתיים בטלפון}

: חפש “תעודת ביטוח רפואי אלקטרונית” ב-Alipay או WeChat. לאחר ההפעלה, תוכל לגשת לטיפול רפואי או לקנות תרופות בבית מרקחת ישירות באמצעות הטלפון, ללא צורך בכרטיס ביטוח רפואי פיזי. בפעם הראשונה נדרש אימות זיהוי פנים, ולאחר מכן בכל שימוש ייווצר קוד דינמי.

\begin{itemize}
  \item \textbf{בירור ביטוח לאומי}: הורד את אפליקציית "Zhangshang 12333" או השתמש ב"מרכז האזרח" ב-Alipay. תוכל לבדוק רשומות תשלומי ביטוח לאומי, יתרת חשבון פנסיה, הגשת בקשה לדמי אבטלה ועוד.
  \item \textbf{מס הכנסה אישי}: הסדר מס הכנסה שנתי בין מרץ ליוני ניתן לבצע באפליקציית "מס הכנסה אישי". זכור לבדוק אם יש ניכויים מיוחדים נוספים שתוכל לדווח עליהם (חינוך ילדים, ריבית משכנתא, תמיכה בהורים מבוגרים ועוד) – אלה יכולים לעזור לך לשלם פחות מס או לקבל החזר מס.
  \item \textbf{זכור את הכללים האלה למניעת הונאות}רשויות אכיפה, בתי משפט, בנקים או נציגי שירות לקוחות לעולם לא יבקשו ממך להעביר כסף ל"חשבון מאובטח".
\end{itemize}

\begin{tcbwarn}{כל מי שמבקש ממך לצלם צילום מסך של קוד התשלום שלך – הוא רמאי.}
\begin{itemize}
  \item אל תלחץ על קישורים לא מוכרים, במיוחד הודעות כמו "פג תוקף ETC" או "כרטיס ביטוח לאומי הושבת".
  \item החזרים על קניות מקוונות מוחזרים באותו מסלול תשלום, אין צורך לספק מספר כרטיס אשראי.
  \item עשה זאת עכשיו: בצע בדיקת אבטחה לאפליקציות התשלום שלך
  \item פתח את Alipay ו-WeChat וודא שתשלום בטביעת אצבע/זיהוי פנים מופעל.
\end{itemize}

\end{tcbwarn}

\begin{tcbpractice}{מצא את פונקציית "מנעול אבטחה" והפעל אותה.}
\begin{enumerate}
  \item בדוק את רשימת "תשלום ללא סיסמה/חיוב אוטומטי" וכבה שירותים שאינך צריך עוד.
  \item בדוק את רשימת "התקנים מחוברים" והסר התקנים לא מוכרים.
  \item הפעל את תעודת הביטוח הרפואי האלקטרונית שלך – בפעם הבאה שתלך לבית חולים לא תצטרך לשאת כרטיס פיזי.
  \item פרק 11: צילום וניהול תמונות
  \item למען האמת, רוב האנשים שמצלמים בטלפון מקבלים תוצאות הרבה מתחת לפוטנציאל האמיתי של המכשיר. זה לא שהטכניקה שלך גרועה, אלא שיש הרבה הגדרות קטנות והרגלים שלא שמת לב אליהם. הפרק הזה לא מדבר באוויר – אלא רק על דברים שתוכל ליישם מיד.
\end{enumerate}

\end{tcbpractice}
